% Converted from README.md
\subsubsection{Callgrind}\label{callgrind}

Callgrind je alat koji generiše listu poziva funkcija korisničkog programa u vidu grafa. U osnovnim podešavanjima sakupljeni podaci obuhvataju broj izvršenih instrukcija, njihov odnos sa linijom u izvršnom kodu, odnos pozivalac/pozvani između funkcija, kao i broj takvih poziva.

Podaci koji se analiziraju nakon završetka rada programa i alata zapisuju se u datoteku \texttt{callgrind.out.\textless{}pid\textgreater{}}, gde \texttt{\textless{}pid\textgreater{}} predstavlja identifikator procesa. Program \texttt{callgrind\_annotate} na osnovu generisanog izveštaja prikazuje listu funkcija. Za grafičku vizuelizaciju preporučeno je koristiti alate kao npr. KCachegrind, koji olakšava navigaciju ukoliko Callgrind napravi veliku količinu podataka.

Program \texttt{callgrind\_control} omogućava interaktivnu kontrolu i nadgledanje programa prilikom izvršavanja. Mogu se dobiti informacije o stanju na steku, a može se i u svakom trenutku generisati profil.

Alat Cachegrind sakuplja podatke, odnosno broji događaje koji se dešavaju direktno u jednoj funkciji. Ovaj mehanizam sakupljanja podataka se naziva ekskluzivnim. Callgrind ovu funkcionalnost proširuje tako što propagira cenu funkcije preko njenih granica. Na primer, ako funkcija \texttt{foo} poziva funkciju \texttt{bar}, cena funkcije \texttt{bar} se dodaje funkciji \texttt{foo}. Kada se ovaj mehanizam primeni na celu funkciju, dobija se slika takozvanih inkluzivnih poziva, gde cena svake funkcije uključuje i cene svih funkcija koje ona poziva, direktno ili indirektno.

Zahvaljujući grafu poziva, može da se odredi, počevši od \texttt{main} funkcije, koja funkcija ima najveću cenu poziva. Odnos pozivalac/pozvani koristan je za profajliranje funkcija koje imaju više poziva iz raznih funkcija. Tada možemo optimizovati program promenom koda u funkciji koja poziva, odnosno smanjenjem broja poziva.

Dodatno može da vrši analizu upotrebe keš memorije i profajliranje grana programa slično kao kod Cachegrind-a. Za takvo profajliranje potrebno je uključiti opcije \texttt{-\/-cache-sim=yes} i \texttt{-\/-branch-sim=yes}. Stoga Callgrind možemo smatrati i proširenjem Cachegrind-a.

\subimport{01_students/}{index.tex}

\subimport{02_prime/}{index.tex}

\subimport{03_pattern_search/}{index.tex}

\subimport{04_sort/}{index.tex}
